% الجزء: sets-functions-relations
% الفصل: sets
% القسم: unions-and-intersections

\documentclass[../../../include/open-logic-section]{subfiles}

\begin{document}

\olfileid[ar]{sfr}{set}{uni}
\olsection{الاتحادات والتقاطعات}

\begin{explain}
سبق في \olref[sfr][set][bas]{sec} تعريف المجموعات بالتجريد، أي بالصورة
$\Setabs{{\OLId{latin-ordinary}{x}}}{\phi({\OLId{latin-ordinary}{x}})}$. والخاصية~$\phi$ المستعملة فيه قد تذكر مجموعات
سبق تعريفها. فإذا كانت ${\OLId{latin-looped}{A}}$ و~${\OLId{latin-looped}{B}}$ مجموعتين، جمعت
$\Setabs{{\OLId{latin-ordinary}{x}}}{{\OLId{latin-ordinary}{x}} \in {\OLId{latin-looped}{A}} \lor {\OLId{latin-ordinary}{x}} \in {\OLId{latin-looped}{B}}}$ كل ما هو من !!{element}s ${\OLId{latin-looped}{A}}$ أو~${\OLId{latin-looped}{B}}$؛
فهي تضم !!{element}s ${\OLId{latin-looped}{A}}$ و~${\OLId{latin-looped}{B}}$ معًا. وفي \olref{fig:union} تصوير ذلك:
فالناحية المظللة تمثل !!{element}s المجموعتين ${\OLId{latin-looped}{A}}$ و~${\OLId{latin-looped}{B}}$ مجتمعة.

\begin{figure}
  \olasset{assets/diagrams/union.tikz}
  \caption{الاتحاد ${\OLId{latin-looped}{A}} \cup {\OLId{latin-looped}{B}}$ يجمع !!{element}s ${\OLId{latin-looped}{A}}$ و~${\OLId{latin-looped}{B}}$ في مجموعة واحدة.}
  \ollabel{fig:union}
\end{figure}

ولكثرة نفع هذه العملية وورودها---أعني جمع المجموعات---نضع لها اسمًا
اصطلاحيًا ورمزًا.
\end{explain}

\begin{defn}[الاتحاد]
\emph{اتحاد} ${\OLId{latin-looped}{A}}$ و${\OLId{latin-looped}{B}}$، ورمزه ${\OLId{latin-looped}{A}} \cup {\OLId{latin-looped}{B}}$، مجموعة الأشياء التي هي من !!{element}s ${\OLId{latin-looped}{A}}$ أو ${\OLId{latin-looped}{B}}$
أو كلتيهما:
\[
{\OLId{latin-looped}{A}} \cup {\OLId{latin-looped}{B}} = \Setabs{{\OLId{latin-ordinary}{x}}}{{\OLId{latin-ordinary}{x}} \in {\OLId{latin-looped}{A}} \lor {\OLId{latin-ordinary}{x}} \in {\OLId{latin-looped}{B}}}
\]
\end{defn}

\begin{ex}
لا أثر لتكرار !!{element}s؛ فإذا اشتركت مجموعتان في !!a{element}
لم يحتج اتحادهما إلى ذكر ذلك !!{element} إلا مرة واحدة، كما في
$\{ {\OLId{latin-ordinary}{a}}, {\OLId{latin-ordinary}{b}}, {\OLId{latin-ordinary}{c}}\} \cup \{ {\OLId{latin-ordinary}{a}}, {\OLMathNumeral{0}}, {\OLMathNumeral{1}}\} = \{{\OLId{latin-ordinary}{a}}, {\OLId{latin-ordinary}{b}}, {\OLId{latin-ordinary}{c}}, {\OLMathNumeral{0}}, {\OLMathNumeral{1}}\}$.

واتحاد المجموعة بإحدى مجموعاتها الجزئية هو المجموعة الأصلية نفسها: $\{{\OLId{latin-ordinary}{a}},
{\OLId{latin-ordinary}{b}}, {\OLId{latin-ordinary}{c}} \} \cup \{{\OLId{latin-ordinary}{a}} \} = \{{\OLId{latin-ordinary}{a}}, {\OLId{latin-ordinary}{b}}, {\OLId{latin-ordinary}{c}}\}$.

وكذلك ضم الخالية إلى مجموعة يبقي المجموعة على حالها: $\{{\OLId{latin-ordinary}{a}},
{\OLId{latin-ordinary}{b}}, {\OLId{latin-ordinary}{c}} \} \cup \emptyset = \{{\OLId{latin-ordinary}{a}}, {\OLId{latin-ordinary}{b}}, {\OLId{latin-ordinary}{c}} \}$.
\end{ex}

\begin{prob}
برهن على أن ${\OLId{latin-looped}{A}} \subseteq {\OLId{latin-looped}{B}}$ يستلزم ${\OLId{latin-looped}{A}} \cup {\OLId{latin-looped}{B}} = {\OLId{latin-looped}{B}}$.
\end{prob}

\begin{explain}
وأما العملية «المزدوجة» للاتحاد فتجمع !!{element}s التي هي من
!!{element}s~${\OLId{latin-looped}{A}}$ ومن !!{element}s~${\OLId{latin-looped}{B}}$ معًا. واسم هذه العملية
\emph{التقاطع}، وصورتها في \olref{fig:intersection}.
\begin{figure}
  \olasset{assets/diagrams/intersection.tikz}
  \caption{التقاطع ${\OLId{latin-looped}{A}} \cap {\OLId{latin-looped}{B}}$ يجمع ما بين المجموعتين من !!{element}s مشتركة.}
  \ollabel{fig:intersection}
\end{figure}
\end{explain}

\begin{defn}[التقاطع]
\emph{تقاطع} ${\OLId{latin-looped}{A}}$ و${\OLId{latin-looped}{B}}$، ورمزه ${\OLId{latin-looped}{A}} \cap {\OLId{latin-looped}{B}}$، مجموعة الأشياء التي هي من
!!{element}s ${\OLId{latin-looped}{A}}$ و~${\OLId{latin-looped}{B}}$ معًا:
\[
{\OLId{latin-looped}{A}} \cap {\OLId{latin-looped}{B}} = \Setabs{{\OLId{latin-ordinary}{x}}}{{\OLId{latin-ordinary}{x}} \in {\OLId{latin-looped}{A}} \land {\OLId{latin-ordinary}{x}} \in {\OLId{latin-looped}{B}}}
\]
فإن خلا التقاطع سميت المجموعتان \emph{متباينتين}؛ أي لا !!{element}s
مشتركة بينهما.
\end{defn}

\begin{ex}
متى انتفت !!{element}s المشتركة بين مجموعتين خلا تقاطعهما، نحو:
$\{ {\OLId{latin-ordinary}{a}}, {\OLId{latin-ordinary}{b}}, {\OLId{latin-ordinary}{c}}\} \cap \{ {\OLMathNumeral{0}}, {\OLMathNumeral{1}}\} = \emptyset$.

ومتى وجدت !!{element}s مشتركة بين مجموعتين، كان تقاطعهما المجموعة
التي تضمها جميعًا ولا تضم غيرها، نحو:
$\{{\OLId{latin-ordinary}{a}}, {\OLId{latin-ordinary}{b}}, {\OLId{latin-ordinary}{c}} \} \cap \{{\OLId{latin-ordinary}{a}}, {\OLId{latin-ordinary}{b}}, {\OLId{latin-ordinary}{d}} \} = \{{\OLId{latin-ordinary}{a}}, {\OLId{latin-ordinary}{b}}\}$.

وتقاطع المجموعة بإحدى مجموعاتها الجزئية هو تلك المجموعة الجزئية:
$\{{\OLId{latin-ordinary}{a}}, {\OLId{latin-ordinary}{b}}, {\OLId{latin-ordinary}{c}}\} \cap \{{\OLId{latin-ordinary}{a}}, {\OLId{latin-ordinary}{b}}\} = \{{\OLId{latin-ordinary}{a}}, {\OLId{latin-ordinary}{b}}\}$.

وأما التقاطع مع الخالية فخالٍ أبدًا: $\{{\OLId{latin-ordinary}{a}}, {\OLId{latin-ordinary}{b}}, {\OLId{latin-ordinary}{c}} \}
\cap \emptyset = \emptyset$.
\end{ex}

\begin{prob}
أقم برهانًا دقيقًا على أن ${\OLId{latin-looped}{A}} \subseteq {\OLId{latin-looped}{B}}$ يستلزم ${\OLId{latin-looped}{A}} \cap {\OLId{latin-looped}{B}} = {\OLId{latin-looped}{A}}$.
\end{prob}

\begin{explain}
وليس الاتحاد والتقاطع مقصورين على مجموعتين. ولتعميمهما وجه موجز:
نجمع أولًا المجموعات المراد اتحادها أو تقاطعها في مجموعة واحدة.
فالاتحاد يجمع كل ما ينتمي إلى !!{element} واحد على الأقل من تلك
المجموعة، وأما التقاطع فيجمع كل ما ينتمي إلى كل !!{element} منها.
\end{explain}

\begin{defn}
لتكن ${\OLId{latin-looped}{A}}$ مجموعة من المجموعات. فاتحادها $\bigcup {\OLId{latin-looped}{A}}$ يجمع
!!{element}s جميع !!{element}s~${\OLId{latin-looped}{A}}$:
\begin{align*}
\bigcup {\OLId{latin-looped}{A}} & = \Setabs{{\OLId{latin-ordinary}{x}}}{{\OLId{latin-ordinary}{x}} \text{ ينتمي إلى !!a{element} من } {\OLId{latin-looped}{A}}},
\text{ أي}\\
& = \Setabs{{\OLId{latin-ordinary}{x}}}{\text{يوجد } {\OLId{latin-looped}{B}} \in {\OLId{latin-looped}{A}}
  \text{ بحيث } {\OLId{latin-ordinary}{x}} \in {\OLId{latin-looped}{B}}}
\end{align*}
\end{defn}

\begin{defn}
لتكن ${\OLId{latin-looped}{A}}$ مجموعة غير خالية من المجموعات. فتقاطعها $\bigcap {\OLId{latin-looped}{A}}$ يجمع ما تشترك
فيه جميع عناصر~${\OLId{latin-looped}{A}}$ من الأشياء:
\begin{align*}
\bigcap {\OLId{latin-looped}{A}} & = \Setabs{{\OLId{latin-ordinary}{x}}}{{\OLId{latin-ordinary}{x}} \text{ ينتمي إلى كل !!{element} من } {\OLId{latin-looped}{A}}},
\text{ أي}\\
 & = \Setabs{{\OLId{latin-ordinary}{x}}}{\text{لكل } {\OLId{latin-looped}{B}} \in {\OLId{latin-looped}{A}}, {\OLId{latin-ordinary}{x}} \in {\OLId{latin-looped}{B}}}
\end{align*}
\end{defn}

\textbf{تنبيه تصحيحي على الأصل.} اشترطنا ألا تخلو عائلة التقاطع؛ لأن
التعريف المعروض لا يعيّن تقاطع العائلة الخالية من غير تعيين مجموعة كلية
تحصر عناصره. ولا يلحق هذا الشرط اتحاد العائلة.

\begin{ex}
لتكن ${\OLId{latin-looped}{A}} = \{ \{ {\OLId{latin-ordinary}{a}}, {\OLId{latin-ordinary}{b}} \}, \{ {\OLId{latin-ordinary}{a}}, {\OLId{latin-ordinary}{d}}, {\OLId{latin-ordinary}{e}} \}, \{ {\OLId{latin-ordinary}{a}}, {\OLId{latin-ordinary}{d}} \} \}$.
فيكون $\bigcup {\OLId{latin-looped}{A}} = \{ {\OLId{latin-ordinary}{a}}, {\OLId{latin-ordinary}{b}}, {\OLId{latin-ordinary}{d}}, {\OLId{latin-ordinary}{e}} \}$ و$\bigcap {\OLId{latin-looped}{A}} = \{ {\OLId{latin-ordinary}{a}} \}$.
\end{ex}
\begin{prob}
	أثبت أنه إذا كانت ${\OLId{latin-looped}{A}}$ مجموعة وكان ${\OLId{latin-looped}{A}} \in {\OLId{latin-looped}{B}}$، فإن
${\OLId{latin-looped}{A}} \subseteq \bigcup {\OLId{latin-looped}{B}}$.
\end{prob}

وكذلك يكون التعريف لمتتالية مجموعات ${\OLId{latin-looped}{A}}_{\OLMathNumeral{1}}$، ${\OLId{latin-looped}{A}}_{\OLMathNumeral{2}}$، \dots
\begin{align*}
\bigcup_{\OLId{latin-ordinary}{i}} {\OLId{latin-looped}{A}}_{\OLId{latin-ordinary}{i}} & = \Setabs{{\OLId{latin-ordinary}{x}}}{{\OLId{latin-ordinary}{x}} \text{ ينتمي إلى إحدى المجموعات } {\OLId{latin-looped}{A}}_{\OLId{latin-ordinary}{i}}}\\
\bigcap_{\OLId{latin-ordinary}{i}} {\OLId{latin-looped}{A}}_{\OLId{latin-ordinary}{i}} & = \Setabs{{\OLId{latin-ordinary}{x}}}{{\OLId{latin-ordinary}{x}} \text{ ينتمي إلى كل } {\OLId{latin-looped}{A}}_{\OLId{latin-ordinary}{i}}}.
\end{align*}

وإن أعطيت \emph{مجموعة فهرسة} ${\OLId{latin-looped}{I}}$، وجعلت مجموعة ${\OLId{latin-looped}{A}}_{\OLId{latin-ordinary}{i}}$ لكل ${\OLId{latin-ordinary}{i}} \in {\OLId{latin-looped}{I}}$،
جاز استعمال اختصار الاتحاد الآتي؛ وأما اختصار التقاطع فنشترط فيه أن تكون
مجموعة الفهرسة غير خالية:
\begin{align*}
	\bigcup_{{\OLId{latin-ordinary}{i}} \in {\OLId{latin-looped}{I}}} {\OLId{latin-looped}{A}}_{\OLId{latin-ordinary}{i}} & = \bigcup \Setabs{{\OLId{latin-looped}{A}}_{\OLId{latin-ordinary}{i}} }{{\OLId{latin-ordinary}{i}} \in {\OLId{latin-looped}{I}}}\\
	\bigcap_{{\OLId{latin-ordinary}{i}} \in {\OLId{latin-looped}{I}}} {\OLId{latin-looped}{A}}_{\OLId{latin-ordinary}{i}} & = \bigcap\Setabs{{\OLId{latin-looped}{A}}_{\OLId{latin-ordinary}{i}}}{{\OLId{latin-ordinary}{i}} \in {\OLId{latin-looped}{I}}}
\end{align*}

ويبقى أن ننظر فيما هو من !!{element}s~${\OLId{latin-looped}{A}}$ وليس في~${\OLId{latin-looped}{B}}$؛ وتمثيله
في \olref{difference}.

\begin{figure}
  \olasset{assets/diagrams/difference.tikz}
  \caption{الفرق ${\OLId{latin-looped}{A}} \setminus {\OLId{latin-looped}{B}}$ يجمع !!{element}s~${\OLId{latin-looped}{A}}$ التي ليست من !!{element}s~${\OLId{latin-looped}{B}}$.}
  \ollabel{difference}
\end{figure}

\begin{defn}[الفرق]
\emph{فرق المجموعتين}~${\OLId{latin-looped}{A}} \setminus {\OLId{latin-looped}{B}}$ هو مجموعة جميع !!{element}s
${\OLId{latin-looped}{A}}$ التي ليست من !!{element}s~${\OLId{latin-looped}{B}}$، أي
\[
{\OLId{latin-looped}{A}}\setminus {\OLId{latin-looped}{B}} = \Setabs{{\OLId{latin-ordinary}{x}}}{{\OLId{latin-ordinary}{x}}\in {\OLId{latin-looped}{A}} \text{ و } {\OLId{latin-ordinary}{x}} \notin {\OLId{latin-looped}{B}}}.
\]
\end{defn}

\begin{prob}
برهن على أن ${\OLId{latin-looped}{A}} \subsetneq {\OLId{latin-looped}{B}}$ يقتضي ${\OLId{latin-looped}{B}} \setminus {\OLId{latin-looped}{A}} \neq \emptyset$.
\end{prob}

\end{document}
